\documentclass[10pt,a4paper]{amsart}
\usepackage[margin=19mm]{geometry}
\usepackage{amsmath,amssymb,mathtools}
\usepackage[scr=rsfs,cal=euler]{mathalfa}
\usepackage{xfrac}
\usepackage[unicode,hidelinks,bookmarks=false]{hyperref}
\newcommand{\ie}{i.e.\ }
\newcommand{\cf}{cf.\ }
\newcommand{\resp}{respectively\ }
\newcommand{\Subsection}{\subsection}
\providecommand{\nobreakdash}{\nobreak\hbox{-}}
\setlength{\emergencystretch}{2em}
\multlinegap0pt
\usepackage{amssymb}
\usepackage{mathtools}
\usepackage[shortlabels]{enumitem}
\newtheorem{lemma}{Lemma}
\newcommand{\BH}{\mathcal B(H)}
\newcommand{\CC}{\mathbb{C}}
\title{Invariant subspaces for operators with spectrum containing the boundary of the numerical range}
\author{Vladimir M\"uller, Yuri Tomilov}
\date{Source: arXiv:2606.22715v1 (CC BY 4.0)}
\begin{document}
\maketitle

\noindent\emph{Source and licence.} Invariant subspaces for operators with spectrum containing the boundary of the numerical range. Version arXiv:2606.22715v1 (CC BY 4.0), distributed by arXiv under the Creative Commons Attribution 4.0 International licence, http://creativecommons.org/licenses/by/4.0/, as recorded in the arXiv licence metadata for this exact version and shown on the abstract page https://arxiv.org/abs/2606.22715v1. Full licence notice: \texttt{licenses/arxiv-2606.22715-CC-BY-4.0.txt}.\par\smallskip

\noindent\textit{Editorial scope.} Counted unit: the article's lemma lem:hull-numerical-range (source0.tex lines 194--202). The article's complete original proof of it is delivered with it (source0.tex lines 204--226, one \texttt{proof} environment of 23 lines). No mathematics is written, repaired or re-derived: the statement and the proof are the article's own lines, in the article's own notation, and no retained line is substituted or re-flowed.  No further source-local context is needed: every cross-reference of every retained line resolves inside the two delivered spans.  Nothing was omitted from any retained span, so the package carries no omission record.  No retained line contains a citation, so the wrapper carries no bibliography and no bibliography entry.  No retained line contains a figure environment, an image-inclusion command or drawing code, so no image is delivered and the package declares no asset dependency.  Editorial formatting only, in addition: an AMS \texttt{amsart} preamble that restates the article's own declarations and macros used by the retained lines, namely the environments \texttt{lemma} on separate counters, together with the standard packages those lines need.  The retained environments print in this document as Lemma 1 in order of appearance, whereas the article prints the same items with the numbering of its own full document.  The complete native source of the article as distributed in the arXiv e-print for this exact version is bundled unchanged as \texttt{sources/203375/source0.tex}.
\begin{lemma}\label{lem:hull-numerical-range}
Let $T\in\BH$ and suppose that $\operatorname{Int} W(T)\ne\emptyset$.
Then
\[
        \partial W(T)\subset\sigma(T)
        \quad\Longleftrightarrow\quad
        \widehat{\sigma(T)}=\overline{W(T)}.
\]
\end{lemma}

\begin{proof}
Put $C=\overline{W(T)}$. Since $W(T)$ is convex and has nonempty
interior, $C$ is a compact convex set with $\partial W(T)=\partial C$.
Moreover $\sigma(T)\subset C$. Suppose first that
$\partial W(T)\subset\sigma(T)$, that is, $\partial C\subset\sigma(T)$.
Then
\[
        \widehat{\partial C}\subset\widehat{\sigma(T)}\subset\widehat C=C .
\]
Since $C$ is compact and convex with nonempty interior,
$\widehat{\partial C}=C$: indeed, $C$ is obtained from $\partial C$ by
filling in the bounded component $\operatorname{Int} C$ of
$\CC\setminus\partial C$. Hence $\widehat{\sigma(T)}=C$.

Conversely, suppose that $\widehat{\sigma(T)}=C$. For every compact plane
set $K$, the set $\widehat K$ is obtained from $K$ by filling in the bounded
components of $\CC\setminus K$.  In particular, one has
$\partial\widehat K\subset K$. Applying this to $K=\sigma(T)$ gives
\[
        \partial W(T)=\partial C=\partial\widehat{\sigma(T)}
        \subset\sigma(T).
\]
\end{proof}

\end{document}
